We explore the algebraic and number-theoretic structure underlying molecular symmetry through the lens of representation theory. We begin by introducing molecular point groups and their matrix representations, focusing on concrete examples of the point groups $C_{2v}$ and $C_{3v}$, which describe the symmetries of water and ammonia molecules, respectively. Then, we construct character tables and decompose the representations into irreducible components. We illustrate how symmetry-adapted linear combinations (SALCs) provide a natural basis for this decomposition.

Building on this representation-theoretic framework, we analyze the eigenvalues and minimal polynomials of symmetry operations. We show that these eigenvalues are roots of unity and hence lie in cyclotomic fields. We also show that for a symmetry operation of order $n$, the minimal polynomial over $\mathbb{Q}$ divides $x^n-1$, hence it factors as a product of cyclotomic polynomials. Motivated by the observation that some character tables contain only integer values despite arising from cyclotomic fields, we investigate the role of Galois automorphisms in governing the structure of character values.

Finally, we extend this analysis to point groups exhibiting fivefold symmetry, where non-integer character values arise. In these cases, we demonstrate that for certain point groups, Galois symmetry shows that relevant character values lie in the real subfield $\mathbb{Q}(\sqrt5)\subset\mathbb{Q}(\zeta_5)$, revealing a deeper connection between molecular symmetry and algebraic number theory. Our central result is that the structure of character tables for molecular point groups is governed by cyclotomic fields and their associated Galois groups.